\chapter{Portfolio Credit}\label{ch:rc:portfolio-credit}

In early 2005 a popular trade sold protection on the equity tranche of the iTraxx
Europe index, the first 3\% of losses on 125 names, and hedged it against market-wide
spread moves by buying protection on the index; the Bank for International Settlements
put its carry at 300 to 400 basis points a year. In May, after the downgrades of General
Motors and Ford (chapter~14), the correlation implied by equity tranche prices fell
sharply, the equity tranche's spread widened, and unwinds pushed the spreads of the
mezzanine tranches down, at times in the opposite direction to the index. Positions
hedged by every model their owners had lost on both legs in the same weeks. This
chapter builds those models: the dependence of defaults through a copula, the loss
distribution of a pool, the prices and risks of tranches, and why a correlation implied
from a tranche is a quoting convention whose moves are a risk of their own. One Quant
Book~2, chapter~24 introduced tranches, the large pool and base correlation; here the
pool is finite, the tranches have a term structure, and the hedges are measured.

\section{Default correlation and copulas}

\begin{definition}[Default correlation]\label{def:rc:portfolio-credit:dc}
The \emph{default correlation}\index{default correlation} of two names over a horizon is
the correlation of their default indicators $\mathbf 1_{\tau_i\le T}$ and
$\mathbf 1_{\tau_j\le T}$:
\[
\rho^{D}_{ij} = \frac{\P(\tau_i\le T,\tau_j\le T)-p_ip_j}{\sqrt{p_i(1-p_i)p_j(1-p_j)}}.
\]
\end{definition}

A copula (One Quant Book~4, chapter~15) joins the marginal distributions of the default
times, each given by a name's survival curve from chapter~13, into a joint
distribution. The market's choice is the simplest.

\begin{definition}[Gaussian copula]\label{def:rc:portfolio-credit:gauss}
In the one-factor \emph{Gaussian copula}\index{Gaussian copula} each name has a latent
variable $X_i = \sqrt\rho\,Z+\sqrt{1-\rho}\,\varepsilon_i$ with $Z,\varepsilon_i$ independent
standard normals, and defaults by $t$ if $X_i < N^{-1}(p_i(t))$, where $p_i(t) = 1-Q_i(t)$.
Each marginal is respected; $\rho$, the correlation of the latent variables (the asset
correlation), sets the dependence.
\end{definition}

The asset correlation is not the \omterm{def:rc:portfolio-credit:dc}{default correlation}. With a five-year default probability
of 4.08\% (a spread of 50 basis points, recovery 40\%), an asset correlation of 25\% gives a
\omterm{def:rc:portfolio-credit:dc}{default correlation} of 7.0\%: rare events of two names coincide less often than their
latent variables co-move.

\section{The one-factor Gaussian copula and the large pool}

Conditional on the common factor $Z = z$ the names are independent, each defaulting by $t$
with probability
\[
p_i(t\mid z) = N\!\left(\frac{N^{-1}(p_i(t))-\sqrt\rho\,z}{\sqrt{1-\rho}}\right).
\]
For a finite pool, the distribution of the number of defaults given $z$ follows by adding
names one at a time: with $P_k$ the probability of $k$ defaults among the first names,
adding name $i$ gives $P_k \leftarrow P_k(1-p_i)+P_{k-1}p_i$, the recursion of Andersen, Sidenius and
Basu (2003). Integrating over $z$ by quadrature gives the unconditional loss distribution
(\cref{fig:rc:portfolio-credit:lossdist}).

\omcode{code/firm/portcredit/firm_portcredit.py}{40}{54}{The recursion over names conditional on the factor, and the unconditional distribution of the number of defaults by Gauss--Hermite quadrature.}\label{lst:rc:portfolio-credit:recursion}

\begin{omfigure}
\begin{tikzpicture}
\begin{axis}[width=11.5cm,height=5.4cm,ymode=log,
  xlabel={number of defaults in five years (125 names)},ylabel={probability},
  tick label style={font=\small},label style={font=\small},
  xmin=0,xmax=40,ymin=1e-6,ymax=1,grid=major,grid style={gray!25},
  legend columns=3,legend style={font=\footnotesize,at={(0.5,-0.25)},anchor=north,draw=none,fill=none,/tikz/every even column/.append style={column sep=8pt}},
  table/col sep=comma]
\addplot[omMeth,very thick] table[x=k,y=r05]{figdata/rates-credit-risk/15-portfolio-credit/loss_dist.csv};
\addplot[omDef,very thick,dashed] table[x=k,y=r25]{figdata/rates-credit-risk/15-portfolio-credit/loss_dist.csv};
\addplot[omThm,very thick,dotted] table[x=k,y=r50]{figdata/rates-credit-risk/15-portfolio-credit/loss_dist.csv};
\legend{$\rho = 5\%$,$\rho = 25\%$,$\rho = 50\%$}
\end{axis}
\end{tikzpicture}

\omcaption{Distribution of the number of defaults in five years in a pool of 125 names
at 50 basis points, for three asset correlations (log scale). The mean is the same, 5.1
defaults; correlation moves probability both to zero defaults and to many. Data: the
chapter's tutorial.}\label{fig:rc:portfolio-credit:lossdist}
\end{omfigure}

\begin{definition}[Large homogeneous pool]\label{def:rc:portfolio-credit:lhp}
The \emph{large homogeneous pool}\index{large homogeneous pool} approximation takes
infinitely many identical names: given $z$ the loss fraction is exactly
$(1-R)\,p(t\mid z)$, so the pool's loss is a monotone function of the factor and its
distribution is available in closed form (the limit Vasicek derived for loan portfolios).
\end{definition}

The large-pool limit of One Quant Book~2 is fast and close for many names: the equity
tranche's five-year expected loss is 51.6\% of its notional in the limit against 49.6\% for
125 names at $\rho = 25\%$. The finite pool matters for the equity tranche and for pools
whose names differ, which is what the recursion is for.

\section{Pricing tranches: recursion, compound and base correlation}

A tranche $[a,d]$ loses $\min(\max(L_t-a,0),d-a)/(d-a)$ of its notional when the pool
has lost $L_t$. Its expected loss at each premium date, $\mathrm{EL}(t_k)$, gives the legs as for
a default swap: the premium is paid on the expected outstanding notional,
$\sum_k\delta_kP(0,t_k)\bigl(1-\tfrac12(\mathrm{EL}(t_{k-1})+\mathrm{EL}(t_k))\bigr)$, and the protection
pays $\sum_kP(0,t_k)\bigl(\mathrm{EL}(t_k)-\mathrm{EL}(t_{k-1})\bigr)$. The equity tranche trades as an upfront
plus 500 basis points running; the others trade at a running spread.

\begin{example}[Synthetic quotes]\label{ex:rc:portfolio-credit:quotes}
Take 125 names at 50 basis points, recovery 40\%, a flat rate of 4\%, and a base-correlation
skew of 15\%, 25\%, 31\%, 36\% and 50\% at detachments 3, 6, 9, 12 and 22\% (illustrative).
The five-year tranches are quoted at an upfront of 37.21\% plus 500 running for 0--3\%, and
at 195, 83.6, 43.6, 17.9 and 3.6 basis points for 3--6, 6--9, 9--12, 12--22 and 22--100\%;
the index is at 49.6.
\end{example}

\begin{definition}[Compound correlation]\label{def:rc:portfolio-credit:compound}
The \emph{compound correlation}\index{compound correlation} of a tranche is the single
correlation at which the one-factor \omterm{def:rc:portfolio-credit:gauss}{Gaussian copula} reprices it, with both its attachment
and detachment losses computed at that correlation.
\end{definition}

The \omterm{def:rc:portfolio-credit:compound}{compound correlations} of the synthetic quotes are 15.0\% for the equity, 12.9\%,
18.4\%, 25.2\% and 50.0\% for the tranches above 6\%, and \emph{two} values for the
mezzanine, 2.1\% and 89.9\%. A mezzanine is long correlation at low correlation (more
correlation spreads losses into it) and short at high correlation (they jump past it), so
its spread is humped in $\rho$ (\cref{fig:rc:portfolio-credit:mezz}). A quote above the hump's
374 basis points, near $\rho = 29\%$, has no \omterm{def:rc:portfolio-credit:compound}{compound correlation} at all. This is why the
market moved to the base correlation of One Quant Book~2, defined on base tranches
$[0,d]$ whose value is monotone in correlation; a mezzanine $[a,d]$ is then priced from two
base correlations, $\mathrm{EL}_{a,d} = (d\,\mathrm{EL}_{0,d}(\rho_d)-a\,\mathrm{EL}_{0,a}(\rho_a))/(d-a)$.

\begin{omfigure}
\begin{tikzpicture}
\begin{axis}[width=11.5cm,height=5.4cm,
  xlabel={compound correlation $\rho$},ylabel={3--6\% par spread (bp)},
  tick label style={font=\small},label style={font=\small},
  xmin=0,xmax=1,ymin=100,ymax=400,grid=major,grid style={gray!25},
  legend columns=2,legend style={font=\footnotesize,at={(0.5,-0.25)},anchor=north,draw=none,fill=none,/tikz/every even column/.append style={column sep=8pt}},
  table/col sep=comma]
\addplot[omDef,very thick] table[x=rho,y=par]{figdata/rates-credit-risk/15-portfolio-credit/mezz.csv};
\addplot[omThm,thick,dashed] table[x=rho,y=quote]{figdata/rates-credit-risk/15-portfolio-credit/mezz.csv};
\legend{model spread,market quote 195 bp}
\end{axis}
\end{tikzpicture}

\omcaption{The 3--6\% tranche's model spread against the \omterm{def:rc:portfolio-credit:compound}{compound correlation}. The
curve is humped, so the quote of 195 basis points is matched at two correlations (2.1\%
and 89.9\%), and a quote above the hump by none. Data: the chapter's
tutorial.}\label{fig:rc:portfolio-credit:mezz}
\end{omfigure}

\begin{definition}[Tranche delta]\label{def:rc:portfolio-credit:delta}
The \emph{tranche delta}\index{tranche delta} is the notional of index protection whose
value changes as much as the tranche's, per unit of tranche notional, when every name's
spread moves together by a small amount, with the base correlations held.
\end{definition}

\begin{example}[Deltas across the structure]\label{ex:rc:portfolio-credit:deltas}
The deltas of the six tranches of \cref{ex:rc:portfolio-credit:quotes} are 17.3, 7.32, 3.28,
1.77, 0.72 and 0.14 (\cref{fig:rc:portfolio-credit:deltas}): EUR~10 million of equity
tranche moves like EUR~173 million of index. Hedging the equity with the mezzanine takes
$17.3/7.32 = 2.36$ times the notional.
\end{example}

\begin{omfigure}
\begin{tikzpicture}
\begin{axis}[width=11.5cm,height=5.2cm,ybar,bar width=16pt,
  symbolic x coords={0-3,3-6,6-9,9-12,12-22,22-100},xtick=data,
  xlabel={tranche (\%)},ylabel={delta (times index)},
  tick label style={font=\small},label style={font=\small},
  ymin=0,ymax=20,grid=major,grid style={gray!25},
  nodes near coords,every node near coord/.append style={font=\footnotesize},
  table/col sep=comma]
\addplot[fill=omDef!70,draw=omDef] table[x=tranche,y=delta]{figdata/rates-credit-risk/15-portfolio-credit/deltas.csv};
\end{axis}
\end{tikzpicture}

\omcaption{\omterm{def:rc:portfolio-credit:delta}{Tranche deltas} to a parallel move of the index, per unit of tranche notional,
with the base correlations held. Leverage is highest at the bottom of the structure: the
equity moves like 17 times its notional of index, the super senior like one seventh. Data:
the chapter's tutorial.}\label{fig:rc:portfolio-credit:deltas}
\end{omfigure}

\begin{definition}[Super senior tranche]\label{def:rc:portfolio-credit:ss}
The \emph{super senior tranche}\index{super senior tranche} is the top of a capital
structure, above the tranche that would earn the highest rating on its own (22--100\% in
the standard European index structure): a very small expected loss on a very large
notional, with most of its risk in the market value of that small loss.
\end{definition}

\begin{remark}[The tail that matters]\label{rem:rc:portfolio-credit:ss}
At an asset correlation of 50\% the chance that the pool's five-year loss exceeds 22\% is
1.88\%: small, on the largest notional of the structure. In 2007--08 the \omterm{def:rc:portfolio-credit:ss}{super senior
tranches} of CDOs backed by mortgage securities, retained by underwriters such as
Citigroup, Merrill Lynch and UBS or insured by AIG, caused billions of dollars of losses,
according to the Financial Crisis Inquiry Commission: holders had treated them as
riskless because their expected loss was. UBS alone wrote down about USD~21.7 billion on \omterm{def:rc:portfolio-credit:ss}{super
senior tranches} it had retained, some USD~50 billion before the crisis and its greatest single source of
loss; the Swiss regulator found the positions had been viewed as VaR-neutral.
\end{remark}

\section{Beyond the Gaussian: fat tails and tail dependence}

\begin{definition}[Student-t copula]\label{def:rc:portfolio-credit:t}
In the \emph{Student-t copula}\index{Student-t copula} with $\nu$ degrees of freedom the
latent variables are $X_i = \sqrt W(\sqrt\rho\,Z+\sqrt{1-\rho}\,\varepsilon_i)$ with $W = \nu/\chi^2_\nu$ common
to all names; a name defaults if $X_i < t_\nu^{-1}(p_i)$. The common scale $W$ makes the
names default together in bad draws: the copula has tail dependence, which the Gaussian
lacks.
\end{definition}

\begin{example}[Tails]\label{ex:rc:portfolio-credit:tails}
At an asset correlation of 25\%, the probability that the pool loses more than 12\% in five
years is 2.24\% under the \omterm{def:rc:portfolio-credit:gauss}{Gaussian copula} and 5.23\% under the \omterm{def:rc:portfolio-credit:t}{Student-t copula} with four
degrees of freedom; for a loss above 24\%, 0.13\% against 1.13\%, 8.6 times more
(\cref{fig:rc:portfolio-credit:tails}). Senior tranches priced with the \omterm{def:rc:portfolio-credit:gauss}{Gaussian copula} at
one correlation look cheap for the same reason the skew exists: the market prices more
tail than the Gaussian gives.
\end{example}

\begin{omfigure}
\begin{tikzpicture}
\begin{axis}[width=11.5cm,height=5.4cm,ymode=log,
  xlabel={pool loss $x$ (\% of notional)},ylabel={$\P(L_5 > x)$},
  tick label style={font=\small},label style={font=\small},
  xmin=0,xmax=30,ymin=1e-4,ymax=1,grid=major,grid style={gray!25},
  legend columns=2,legend style={font=\footnotesize,at={(0.5,-0.25)},anchor=north,draw=none,fill=none,/tikz/every even column/.append style={column sep=8pt}},
  table/col sep=comma]
\addplot[omDef,very thick] table[x=x,y=gauss]{figdata/rates-credit-risk/15-portfolio-credit/tails.csv};
\addplot[omThm,very thick,dashed] table[x=x,y=t4]{figdata/rates-credit-risk/15-portfolio-credit/tails.csv};
\legend{Gaussian copula,Student-t copula ($\nu = 4$)}
\end{axis}
\end{tikzpicture}

\omcaption{Probability that the five-year pool loss exceeds $x$ (log scale), same
marginals and asset correlation 25\%, simulated with 400\,000 draws of the factors. The
\omterm{def:rc:portfolio-credit:t}{Student-t copula} puts several times more probability on large losses. Data: the chapter's
tutorial.}\label{fig:rc:portfolio-credit:tails}
\end{omfigure}

\section{May 2005 and 2008}

\begin{dated}{2026-09}{May 2005}\label{dat:rc:portfolio-credit:may2005}
According to the Bank for International Settlements' review of June 2005, selling
protection on the iTraxx Europe equity tranche hedged with index protection earned 300 to
400 basis points in early 2005; in May the correlation implied by equity tranche prices
fell sharply, equity spreads widened (faster after the downgrades of General Motors and
Ford), unwinds pushed mezzanine spreads down, and at times the mezzanine moved opposite
to the index.
\end{dated}

The two episodes stress different parts of the model. In May 2005 the risk was
idiosyncratic: one or two names blowing out hit the equity tranche, which absorbs the first
losses, far more than an index-sized hedge, and the correlation implied by equity prices
fell; the delta, computed for a parallel move with correlation held, hedged neither. In
2008 the risk was systematic: the factor moved, the senior tranches that the \omterm{def:rc:portfolio-credit:gauss}{Gaussian
copula} had priced as nearly riskless lost value, and the correlations implied by senior
tranches rose towards one.

\section{Tutorial: tranches by recursion}\label{tut:rc:portfolio-credit:tutorial}

\begin{tutorial}
\textbf{Goal.} Price a tranche structure on a 125-name pool, imply \omterm{def:rc:portfolio-credit:compound}{compound
correlations}, compare copulas and compute deltas. \textbf{End state:} the numbers of
\cref{ex:rc:portfolio-credit:quotes,ex:rc:portfolio-credit:deltas,ex:rc:portfolio-credit:tails} and
\cref{fig:rc:portfolio-credit:lossdist,fig:rc:portfolio-credit:mezz,fig:rc:portfolio-credit:deltas,fig:rc:portfolio-credit:tails}.
\begin{tutsteps}
\item \textbf{Loss distribution}: \texttt{loss\_distribution(p, rho)} for three
correlations.
\item \textbf{Quotes}: \texttt{quotes()} from the base-correlation skew.
\item \textbf{Compound}: \texttt{compound\_table()}; plot \texttt{mezz\_curve()}.
\item \textbf{Tails and deltas}: \texttt{tail\_table()}, \texttt{deltas()};
\texttt{fig\_rc\_portcredit.py} writes the charts.
\end{tutsteps}
\textbf{What to change next.} Give ten names a spread of 300 basis points and the rest 40,
at the same index level, and compare the equity tranche's price; raise the mezzanine
quote to 400 and watch the \omterm{def:rc:portfolio-credit:compound}{compound correlation} disappear.
\end{tutorial}

\section{Build: portfolio credit}\label{bld:rc:portfolio-credit:portcredit}

\begin{build}
\bfield{Purpose} The firm's pool models: loss distributions of credit portfolios for
tranches, for the counterparty portfolios of Part~III, and for credit capital in Part~IV.
\bfield{Interface} \texttt{Pool(hazards, recovery)}; \texttt{conditional\_pd},
\texttt{default\_count\_given\_z}, \texttt{loss\_distribution}; \texttt{expected\_tranche\_losses},
\texttt{base\_el}, \texttt{legs\_from\_el}, \texttt{tranche\_price};
\texttt{compound\_correlations}; \texttt{t\_cdf}, \texttt{t\_inv}, \texttt{simulate\_losses};
\texttt{default\_correlation}.
\bfield{Rules} Flat hazards per name, equal notionals, one recovery; quarterly premiums on
the expected outstanding notional; 64-point Gauss--Hermite quadrature; Book~2's
\texttt{firm\_tranche} is the large-pool limit.
\bfield{Acceptance tests} \texttt{code/firm/portcredit/tests/}: the recursion without
correlation is binomial; a large pool approaches Book~2's limit; base and compound agree
on a flat skew; the index does not depend on correlation; Student-t quantiles; simulated
mean losses.
\bfield{Stretch} Unequal notionals and recoveries (losses on a grid); stochastic recovery;
the random-factor-loading and Student-t pricers by quadrature; bespoke pools mapped to
index base correlations.
\end{build}

\begin{omsources}
\item D.\ X.\ Li, ``On default correlation: a copula function approach'', \emph{Journal of
Fixed Income} 9(4), 2000, 43--54.
\item O.\ A.\ Vasicek, \emph{Finance, Economics and Mathematics}, Wiley, 2015, chapter
``Loan portfolio value''.
\item J.\ Hull and A.\ White, ``Valuation of a CDO and an $n$-th to default CDS without
Monte Carlo simulation'', \emph{Journal of Derivatives} 12(2), 2004, 8--23.
\item Bank for International Settlements, \emph{Quarterly Review}, June 2005.
\item Financial Crisis Inquiry Commission, \emph{The Financial Crisis Inquiry Report},
2011.
\item L.\ Andersen, J.\ Sidenius and S.\ Basu, ``All your hedges in one basket'', \emph{Risk},
November 2003.
\item Swiss Federal Banking Commission, \emph{Subprime crisis: SFBC investigation into the causes of the
write-downs of UBS AG}, 30 September 2008.
\end{omsources}

\section{Exercises}

\begin{exercise}[$\star$]\label{exo:rc:portfolio-credit:1}
A pool of 125 names each has a five-year default probability of 4.08\%. How many defaults
are expected, and does the expectation depend on correlation?
\end{exercise}

\begin{exercise}[$\star$]\label{exo:rc:portfolio-credit:2}
Which tranches are long correlation and which short, and why is the mezzanine both?
\end{exercise}

\begin{exercise}[$\star$]\label{exo:rc:portfolio-credit:3}
With a recovery of 40\%, how many defaults among 125 names wipe out the 0--3\% tranche,
and how many reach the super senior at 22\%?
\end{exercise}

\begin{exercise}[$\star\star$]\label{exo:rc:portfolio-credit:4}
Explain why an asset correlation of 25\% gives a \omterm{def:rc:portfolio-credit:dc}{default correlation} of only 7\%.
\end{exercise}

\begin{exercise}[$\star\star$]\label{exo:rc:portfolio-credit:5}
Using the deltas of \cref{ex:rc:portfolio-credit:deltas}, how much index protection hedges
EUR~50 million of sold 6--9\% protection?
\end{exercise}

\begin{exercise}[$\star\star$]\label{exo:rc:portfolio-credit:6}
Why does the market quote base correlations rather than \omterm{def:rc:portfolio-credit:compound}{compound correlations}, and
what is wrong with base correlation?
\end{exercise}

\begin{exercise}[$\star\star\star$]\label{exo:rc:portfolio-credit:7}
\emph{Coding.} Compute the probability that the pool loses more than 24\% in five years under
both copulas at an asset correlation of 25\%, and the ratio.
\end{exercise}

\begin{exercise}[$\star\star\star$]\label{exo:rc:portfolio-credit:8}
\emph{Find the flaw.} ``Our super senior position has an expected loss of a few basis points
and an AAA rating. It needs no hedge and almost no capital.''
\end{exercise}

\section{Problem: Long Equity, Short Mezzanine}

\begin{problem}[{Weekend problem --- the correlation trade of May 2005}]\label{pb:rc:portfolio-credit:1}
A fund sells EUR~10 million of five-year 0--3\% protection at the quotes of
\cref{ex:rc:portfolio-credit:quotes} (upfront 37.21\% plus 500 running) and buys 3--6\%
protection at 195 basis points in the delta-neutral amount. Then one name's spread jumps
from 50 to 1\,000 basis points, and the equity base correlation falls from 15\% to 12\%,
the other base correlations unchanged.

\medskip\noindent\textbf{Part I --- The position.}
\begin{enumerate}
\item Give the deltas of the two tranches and the mezzanine notional.
\item What does the fund receive at inception, and what does it earn while nothing
happens?
\item What is its exposure to a parallel widening, to one name defaulting, and to
correlation?
\item Why was the trade called long correlation?
\item Which risks did the delta measure not include?
\end{enumerate}

\medskip\noindent\textbf{Part II --- One name blows out.}
\begin{enumerate}[resume]
\item Give the equity upfront and mezzanine spread after the jump.
\item Give the P\&L of each leg and the total.
\item Why does the hedge not cover the equity's loss?
\item What would one default of that name (recovery 40\%) cost the equity tranche
directly, in notional?
\item Would an index hedge have done better?
\end{enumerate}

\medskip\noindent\textbf{Part III --- Correlation falls.}
\begin{enumerate}[resume]
\item Give the P\&L of each leg when only the equity base correlation falls to 12\%.
\item Why does the mezzanine spread fall?
\item Give both legs' P\&L and the total when both happen.
\item Which leg did the fund expect to protect it, and what happened to it?
\item How does the unwinding of similar positions move the prices further?
\end{enumerate}

\medskip\noindent\textbf{Part IV --- Judgement.}
\begin{enumerate}[resume]
\item What is the correlation that the fund was long, in the model and in the market?
\item How would you limit such a book?
\item Why did the same models give no warning of 2008's senior losses?
\item State the \emph{named result}: the two legs' P\&L and the total when one name blows
out and equity correlation falls.
\item In one sentence: what does a delta-hedged tranche position still bet on?
\end{enumerate}
\end{problem}

\section{Interview questions}

\begin{interviewq}[$\star$ \iqroles{researcher, trader}]\label{iq:rc:portfolio-credit:1}
Write down the one-factor \omterm{def:rc:portfolio-credit:gauss}{Gaussian copula} and the conditional default probability.
\end{interviewq}

\begin{interviewq}[$\star\star$ \iqroles{researcher, developer}]\label{iq:rc:portfolio-credit:2}
How do you compute the loss distribution of a finite pool of heterogeneous names
without Monte Carlo?
\end{interviewq}

\begin{interviewq}[$\star\star$ \iqroles{trader}]\label{iq:rc:portfolio-credit:3}
Compound versus base correlation: what is each, and what goes wrong with each?
\end{interviewq}

\begin{interviewq}[$\star\star$ \iqroles{risk}]\label{iq:rc:portfolio-credit:4}
How would you risk-manage a book of \omterm{def:rc:portfolio-credit:ss}{super senior tranches}?
\end{interviewq}

\begin{interviewq}[$\star\star\star$ \iqroles{researcher}]\label{iq:rc:portfolio-credit:5}
What is tail dependence, which copulas have it, and why does it matter for senior
tranches?
\end{interviewq}

\begin{interviewq}[$\star\star\star$ \iqroles{trader, risk}]\label{iq:rc:portfolio-credit:6}
Explain how a delta-hedged long equity, short mezzanine position lost money in May 2005.
\end{interviewq}
